% 题证摘录；来源与整理边界见相同编号分题文件。

\paragraph{53.12 Riemann--Hurwitz: the proof discussion preceding the numbered statement.}
Let $k$ be a field. Let $f:X\to Y$ be a morphism of smooth curves over $k$. Then we obtain a canonical exact sequence
\[f^*\Omega_{Y/k}\xrightarrow{\mathrm df}\Omega_{X/k}\longrightarrow\Omega_{X/Y}\longrightarrow0\]
by Morphisms, Lemma 29.33.9. Since $X$ and $Y$ are smooth, the sheaves $\Omega_{X/k}$ and $\Omega_{Y/k}$ are invertible modules, see Morphisms, Lemma 29.35.12. Assume the first map is nonzero, i.e., assume $f$ is generically \'etale, see Lemma 53.12.1. Let $R\subset X$ be the closed subscheme cut out by the different $\mathfrak D_f$ of $f$. By Discriminants, Lemma 49.12.6 this is the same as the vanishing locus of $\mathrm df$, it is an effective Cartier divisor, and we get
\[f^*\Omega_{Y/k}\otimes_{\mathcal O_X}\mathcal O_X(R)=\Omega_{X/k}.\]
In particular, if $X,Y$ are projective with $k=H^0(Y,\mathcal O_Y)=H^0(X,\mathcal O_X)$ and $X,Y$ have genus $g_X,g_Y$, then we get the Riemann--Hurwitz formula
\begin{align*}
2g_X-2&=\deg(\Omega_{X/k})\\
&=\deg\bigl(f^*\Omega_{Y/k}\otimes_{\mathcal O_X}\mathcal O_X(R)\bigr)\\
&=\deg(f)\deg(\Omega_{Y/k})+\deg(R)\\
&=\deg(f)(2g_Y-2)+\deg(R).
\end{align*}
The first and last equality by Lemma 53.8.4. The second equality by the isomorphism of invertible sheaves given above. The third equality by additivity of degrees (Varieties, Lemma 33.44.7), the formula for the degree of a pullback (Varieties, Lemma 33.44.11), and finally the formula for the degree of $\mathcal O_X(R)$ (Varieties, Lemma 33.44.9).

To use the Riemann--Hurwitz formula we need to compute $\deg(R)=\dim_k\Gamma(R,\mathcal O_R)$. By the structure of zero dimensional schemes over $k$ (see for example Varieties, Lemma 33.20.2), we see that $R$ is a finite disjoint union of spectra of Artinian local rings $R=\coprod_{x\in R}\Spec(\mathcal O_{R,x})$ with each $\mathcal O_{R,x}$ of finite dimension over $k$. Thus
\[\deg(R)=\sum_{x\in R}\dim_k\mathcal O_{R,x}=\sum_{x\in R}d_x[\kappa(x):k]\]
with
\[d_x=\operatorname{length}_{\mathcal O_{R,x}}\mathcal O_{R,x}=\operatorname{length}_{\mathcal O_{X,x}}\mathcal O_{R,x}\]
the multiplicity of $x$ in $R$ (see Algebra, Lemma 10.52.12). Let $x\in X$ be a closed point with image $y\in Y$. Looking at stalks we obtain an exact sequence
\[\Omega_{Y/k,y}\to\Omega_{X/k,x}\to\Omega_{X/Y,x}\to0.\]
Choosing local generators $\eta_x$ and $\eta_y$ of the (free rank $1$) modules $\Omega_{X/k,x}$ and $\Omega_{Y/k,y}$ we see that $\eta_y\mapsto h\eta_x$ for some nonzero $h\in\mathcal O_{X,x}$. By the exact sequence we see that $\Omega_{X/Y,x}\cong\mathcal O_{X,x}/h\mathcal O_{X,x}$ as $\mathcal O_{X,x}$-modules. Since the divisor $R$ is cut out by $h$ (see above) we have $\mathcal O_{R,x}=\mathcal O_{X,x}/h\mathcal O_{X,x}$. Thus we find the following equalities
\begin{align*}
d_x&=\operatorname{length}_{\mathcal O_{X,x}}(\mathcal O_{R,x})\\
&=\operatorname{length}_{\mathcal O_{X,x}}(\mathcal O_{X,x}/h\mathcal O_{X,x})\\
&=\operatorname{length}_{\mathcal O_{X,x}}(\Omega_{X/Y,x})\\
&=\operatorname{ord}_{\mathcal O_{X,x}}(h)\\
&=\operatorname{ord}_{\mathcal O_{X,x}}(\text{``}\eta_y/\eta_x\text{''}).
\end{align*}
The first equality by our definition of $d_x$. The second and third we saw above. The fourth equality is the definition of $\operatorname{ord}$, see Algebra, Definition 10.121.2. Note that since $\mathcal O_{X,x}$ is a discrete valuation ring, the integer $\operatorname{ord}_{\mathcal O_{X,x}}(h)$ just the valuation of $h$. The fifth equality is a mnemonic.

Here is a case where one can ``calculate'' the multiplicity $d_x$ in terms of other invariants. Namely, if $\kappa(x)$ is separable over $k$, then we may choose $\eta_x=\mathrm ds$ and $\eta_y=\mathrm dt$ where $s$ and $t$ are uniformizers in $\mathcal O_{X,x}$ and $\mathcal O_{Y,y}$ (Lemma 53.12.3). Then $t\mapsto us^{e_x}$ for some unit $u\in\mathcal O_{X,x}$ where $e_x$ is the ramification index of the extension $\mathcal O_{Y,y}\subset\mathcal O_{X,x}$. Hence we get
\[\eta_y=\mathrm dt=\mathrm d(us^{e_x})=es^{e_x-1}u\,\mathrm ds+s^{e_x}\mathrm du.\]
Writing $\mathrm du=w\,\mathrm ds$ for some $w\in\mathcal O_{X,x}$ we see that
\[\text{``}\eta_y/\eta_x\text{''}=es^{e_x-1}u+s^{e_x}w=(e_xu+sw)s^{e_x-1}.\]
We conclude that the order of vanishing of this is $e_x-1$ unless the characteristic of $\kappa(x)$ is $p>0$ and $p$ divides $e_x$ in which case the order of vanishing is $>e_x-1$.

Combining all of the above we find that if $k$ has characteristic zero, then
\[2g_X-2=(2g_Y-2)\deg(f)+\sum_{x\in X}(e_x-1)[\kappa(x):k]\]
where $e_x$ is the ramification index of $\mathcal O_{X,x}$ over $\mathcal O_{Y,f(x)}$.

\begin{lemma}[53.12.2; Tag 0C1D]
Let $f:X\to Y$ be a morphism of smooth proper curves over a field $k$ which satisfies the equivalent conditions of Lemma 53.12.1. If $k=H^0(Y,\mathcal O_Y)=H^0(X,\mathcal O_X)$ and $X$ and $Y$ have genus $g_X$ and $g_Y$, then
\[2g_X-2=(2g_Y-2)\deg(f)+\deg(R)\]
where $R\subset X$ is the effective Cartier divisor cut out by the different of $f$.
\end{lemma}
\begin{proof}
See discussion above; we used Discriminants, Lemma 49.12.6, Lemma 53.8.4, and Varieties, Lemmas 33.44.7 and 33.44.11.
\end{proof}
\begin{lemma}[53.12.3; Tag 0C1E]
Let $X\to\Spec(k)$ be smooth of relative dimension $1$ at a closed point $x\in X$. If $\kappa(x)$ is separable over $k$, then for any uniformizer $s$ in the discrete valuation ring $\mathcal O_{X,x}$ the element $\mathrm ds$ freely generates $\Omega_{X/k,x}$ over $\mathcal O_{X,x}$.
\end{lemma}
\begin{proof}
The ring $\mathcal O_{X,x}$ is a discrete valuation ring by Algebra, Lemma 10.140.3. Since $x$ is closed $\kappa(x)$ is finite over $k$. Hence if $\kappa(x)/k$ is separable, then any uniformizer $s$ maps to a nonzero element of $\Omega_{X/k,x}\otimes_{\mathcal O_{X,x}}\kappa(x)$ by Algebra, Lemma 10.140.4. Since $\Omega_{X/k,x}$ is free of rank $1$ over $\mathcal O_{X,x}$ the result follows.
\end{proof}
